\Needspace{6\baselineskip}
\begin{center}\hypertarget{native-262-article189}{}\pdfbookmark[2]{189.}{article-189}189.\end{center}
This method differs from the preceding one only in that it fits two geocentric
positions exactly and the remaining positions as closely as possible. The
other method gives no observation preference over the others and distributes
the errors among them all as far as possible\rlap{\raisebox{1.4ex}{\rule{.6em}{.35pt}}}.
The method of the preceding article is therefore inferior to the earlier one
only insofar as the errors in the remaining positions can be appreciably
reduced by allowing the positions $L,L'$ to take on some of the error.
Usually, however, a suitable choice of observations $L,L'$ can easily keep
this difference from becoming important. Care must \SourceLeaf{263}be taken to
choose observations for $L,L'$ that are both exceptionally precise and such
that the elements derived from them and the distances are not unduly affected
by small changes in the geocentric positions themselves. It would therefore
be unwise to choose observations separated by a short interval of time or
observations corresponding to heliocentric positions that are nearly opposite
or nearly coincident.
\begin{center}\rule{3cm}{.4pt}\end{center}
